\chapter{Correlation Metric}
Our metric captures two aspects simultaneously: (i) how different the query-local attribute distribution is from the global dataset distribution, and (ii) whether the queried filter range is locally enriched or depleted.

Formally, let \(X = \{x_1,\dots,x_N\} \subseteq \mathbb{R}^d\) be the dataset of embedding vectors, and let each vector \(x_i\) be associated with a scalar attribute \(a_i\). 
Given a query vector \(q \in \mathbb{R}^d\) and a filter range \([L,U]\), we seek to determine whether values in \([L,U]\) are more likely in the neighborhood of \(q\) than in the dataset as a whole.

\section{Precomputed statistics}
We precompute three components: 

\subsection{Cluster centroids.}
Let \(\ell_i \in \{1,\dots,k\}\) denote the cluster assignment of vector \(x_i\). 
For each cluster \(c\), define
\[
C_c = \{ i \mid \ell_i = c \}.
\]
The centroid is reconstructed as
\[
c_c = \frac{1}{|C_c|} \sum_{i \in C_c} x_i.
\]

\subsection{Attribute distributions.}
We estimate both global and cluster-local attribute distributions using histogram-based density estimation over the scalar attribute \(a_i\).

\textbf{Bin construction.}
The number of histogram bins $M$ is determined adaptively from the data.
We begin by estimating the bin width using the Freedman--Diaconis rule,
\[
h = \frac{2 \cdot \mathrm{IQR}}{n^{1/3}},
\]
where $\mathrm{IQR}$ is the interquartile range of the attribute values and $n$
is the number of samples. Given the data range $[v_{\min}, v_{\max}]$, this
yields the bin-count estimate
\[
\widehat{M} = \left\lceil \frac{v_{\max} - v_{\min}}{h} \right\rceil.
\]
When the interquartile range is too small (i.e., the data is nearly constant),
this estimate becomes unreliable, and we instead fall back to a heuristic based
on a target number of samples per bin.

To keep the number of bins reasonable, we force it to stay within a fixed range:
\[
M = \min\!\bigl(M_{\max},\, \max(M_{\min},\, \widehat{M})\bigr).
\]
This way, the histogram is neither too rough (losing detail) nor too detailed
(with many empty bins that make the results unreliable).

In our implementation, we fix $M_{\min} = 20$ and $M_{\max} = 200$. Since the
dataset characteristics are known, these bounds are
sufficient and yield stable results in practice.

More generally, these bounds can be defined in a data-dependent manner. For
example, letting $u$ denote the number of distinct attribute values, one may use
\[
M_{\min} = \min\!\left(u,\, \max\!\left(5, \left\lfloor n^{1/3} \right\rfloor \right)\right),
\qquad
M_{\max} = \min\!\left(u,\, \left\lfloor \frac{n}{m_{\min}} \right\rfloor \right),
\]
where $m_{\min}$ is the minimum desired number of samples per bin. Such choices
adapt the binning resolution to both dataset size and attribute cardinality.

\textbf{Global distribution.}
Given bin edges \(b_0 < b_1 < \dots < b_M\), we first compute a histogram over
the full dataset:
\[
H_{\mathrm{global}}(j)
=
\left|
\{\, i \mid a_i \in [b_j, b_{j+1}) \,\}
\right|,
\qquad j=0,\dots,M-1.
\]

We then convert this histogram into a probability distribution. To avoid zero
probabilities, we apply additive smoothing by adding a small constant
\(\varepsilon > 0\) to each bin:
\[
P_{\mathrm{global}}(j)
=
\frac{H_{\mathrm{global}}(j) + \varepsilon}
{\sum_{k=0}^{M-1} \bigl(H_{\mathrm{global}}(k) + \varepsilon\bigr)}.
\]
This ensures that no bin has zero probability.
Without smoothing, bins with zero counts would lead to undefined or unstable
values in downstream computations (e.g., divisions or logarithms in divergence
measures).

After smoothing, the histogram no longer sums to one, so we renormalize it.
We simply rescale all values so that they form a valid probability
distribution again:
\[
\sum_j P_{\mathrm{global}}(j) = 1.
\]

The corresponding cumulative distribution function is
\[
\mathrm{CDF}_{\mathrm{global}}(j)
=
\sum_{k=0}^{j} P_{\mathrm{global}}(k).
\]

\textbf{Cluster-local distributions.}
For each cluster \(c\), we first compute its radius
\[
r_c = \max_{i \in C_c} \|x_i - c_c\|_2.
\]
We then introduce a \textbf{scaling parameter} \(\lambda \ge 1\) to slightly enlarge the
cluster and define the set of selected points
\[
C'_c = \{\, i \in C_c \mid \|x_i - c_c\|_2 \le \lambda r_c \,\}.
\]
This allows us to include points near the boundary of the cluster, making the
estimate more robust. The choice and effect of \(\lambda\) are discussed in
Section~\ref{subsec:radius-expansion}.

Using these points, we estimate the cluster-level attribute distribution by
constructing a histogram over the same bins:
\[
H_c(j)
=
\left|
\{\, i \in C'_c \mid a_i \in [b_j, b_{j+1}) \,\}
\right|,
\qquad j=0,\dots,M-1,
\]
which corresponds to the empirical distribution
\[
P_c(a) = \frac{1}{|C'_c|} \sum_{i \in C'_c} \mathbf{1}[a_i \in \text{bin}(a)].
\]
We then convert it into a probability distribution by applying additive
smoothing and renormalization, as in the global case:
\[
P_c(j)
=
\frac{H_c(j) + \varepsilon}
{\sum_{k=0}^{M-1} \bigl(H_c(k) + \varepsilon\bigr)}.
\]
Finally, the cumulative distribution function is computed analogously:
\[
\mathrm{CDF}_c(j)
=
\sum_{k=0}^{j} P_c(k).
\]

\paragraph{Query-conditional distribution.}
We consider the \emph{query-conditional attribute distribution}
\(P(a \mid x)\), which describes the distribution of attribute values in the
neighborhood of a point \(x \in \mathbb{R}^d\). Intuitively, it captures how
attribute values are distributed among vectors that are similar to \(x\) in the
embedding space.

Estimating \(P(a \mid x)\) directly is challenging, as it requires defining a
local neighborhood around each query and computing a distribution from a
potentially small number of samples. Instead, we approximate the mapping
\[
x \mapsto P(a \mid x)
\]
by assigning a single distribution \(P_c(a)\) to each cluster:
\[
x \mapsto P_c(a), \quad \text{for } x \in R_c,
\]
where \(R_c\) denotes the region associated with cluster \(c\). This avoids
estimating a separate distribution for every query location, which would be
both computationally expensive and statistically unreliable.

\paragraph{Distributions and estimation.}
We distinguish between three distributions. 

The conditional distribution
\(P(a \mid x)\) is defined at the point level. 

At the cluster level, we
introduce an ideal (population-level) distribution \(P_c(a)\), which we define
as the average of the conditional distributions over the region \(R_c\):
\[
P_c(a) = \mathbb{E}_{x' \sim \mathrm{Unif}(R_c)}[P(a \mid x')].
\]
This provides a clean abstraction of the attribute distribution associated
with the cluster.

In practice, this distribution is not directly observable and is instead
estimated empirically using a histogram over the points in the cluster:
\[
\hat{P}_c(a) = \frac{1}{n_c} \sum_{i \in C'_c} \mathbf{1}[a_i \in \text{bin}(a)].
\]

\subsection{Query-to-cluster mapping.}
Given a query \(q \in \mathbb{R}^d\), we associate it with a representative
region of the dataset by assigning it to its nearest centroid. To this end, we
construct an index over the set of centroids \(\{c_1, \dots, c_k\}\) and define
\[
c^*(q) = \arg\min_c \|q - c_c\|_2.
\]
This induces a partition of the vector space into Voronoi regions and assigns
each query to the cluster whose centroid is closest to it.

To estimate the attribute distribution in the vicinity of \(q\), we use the
empirical distribution of its assigned cluster, given by the histogram-based
estimate computed in the previous section. 
\[
P(a \mid q) \approx P_{c^*(q)}(a),
\]
where \(P_{c^*(q)}\) denotes the attribute distribution associated with the
cluster \(c^*(q)\). 

\section{Approximation quality}\label{sec:approximation-quality}
We now analyze how well the cluster-based approximation
\[
P(a \mid q) \approx P_c(a)
\]
captures the true query-conditional distribution.

\paragraph{Smoothness assumption}
We assume that attribute distributions are correlated with the underlying
vector space geometry, meaning that nearby points exhibit similar attribute
behavior. 

Formally, we assume a \emph{Lipschitz continuity} condition: there exists a
constant \(L > 0\) such that
\[
D(P(a \mid x), P(a \mid x')) \le L \|x - x'\|
\]
for all \(x, x' \in \mathbb{R}^d\), where \(D(\cdot,\cdot)\) measures the
difference between two distributions.

This means that the change in the attribute distribution is proportional to the
distance between points: small movements in the embedding space lead to small
changes in the distribution, while large differences can only occur between
points that are far apart.

\subsection{Bounding the approximation error}

For a point \(x \in R_c\), the approximation error is
\[
D(P(a \mid x), P_c(a)),
\]
where \(D(\cdot,\cdot)\) is a divergence between distributions.

Recall that we define the ideal cluster-level distribution as the average of
the conditional distributions over the region \(R_c\), i.e., \(P_c(a)\) is the
average of \(\{P(a \mid x') : x' \in R_c\}\).

We now bound the deviation between \(P(a \mid x)\) and this average. Since
\(P_c(a)\) is an average over these distributions, its distance to
\(P(a \mid x)\) can be bounded by the largest deviation within the cluster:
\[
D(P(a \mid x), P_c(a))
\le
\sup_{x' \in R_c} D(P(a \mid x), P(a \mid x')).
\]

Applying the smoothness assumption gives:
\[
D(P(a \mid x), P_c(a))
\le
\sup_{x' \in R_c} D(P(a \mid x), P(a \mid x'))
\]
\[
D(P(a \mid x), P_c(a))
\le
\sup_{x' \in R_c} L \|x - x'\|
\]
\[
D(P(a \mid x), P_c(a))
\le
L \sup_{x' \in R_c} \|x - x'\|.
\]

Taking the supremum over all \(x \in R_c\) yields:
\[
\sup_{x \in R_c} D(P(a \mid x), P_c(a))
\le
L \sup_{x,x' \in R_c} \|x - x'\|.
\]

\paragraph{Cluster diameter.}
We define the diameter of a cluster \(R_c\) as
\[
\mathrm{diam}(R_c) = \sup_{x,x' \in R_c} \|x - x'\|.
\]
Substituting this into the previous bound yields
\[
\sup_{x \in R_c} D(P(a \mid x), P_c(a))
\le
L \, \mathrm{diam}(R_c).
\]


\subsection{Error decomposition}

The proposed approximation introduces two main sources of error. 

The first is the quantization (approximation) error,
\[
P(a \mid x) \;\;\text{vs}\;\; P_c(a),
\]
which arises from replacing the pointwise conditional distribution with a
cluster-level approximation. 

The second is the estimation error,
\[
P_c(a) \;\;\text{vs}\;\; \hat{P}_c(a),
\]
which is introduced when estimating the cluster-level distribution from a
finite number of samples using a histogram.

\subsubsection{Quantization error (Geometric Discretization).}
For a query assigned to cluster \(c\), we define the worst-case approximation
error as
\[
\varepsilon_{\mathrm{quant}}(q)
=
\sup_{x \in R_c} D(P(a \mid x), P_c(a)),
\]
where \(D(\cdot,\cdot)\) is a divergence between distributions (e.g., total
variation or Jensen--Shannon). This measures the error introduced by replacing
the pointwise distribution \(P(a \mid x)\) with a single cluster-level
distribution.

From the previous bound, this error satisfies
\[
\varepsilon_{\mathrm{quant}}(q)
\le
L \, \mathrm{diam}(R_c),
\]
where \(\mathrm{diam}(R_c)\) is the cluster diameter and \(L\) is the smoothness
constant. Thus, the approximation error increases with the size of the cluster
and with how rapidly the distribution changes in the embedding space.

However, this bound is worst-case and does not fully reflect practical behavior.
In particular, the error also depends on how the geometry of the cluster aligns
with the structure of \(P(a \mid x)\).

\paragraph{Refinements beyond diameter.}
Three factors are particularly important:

\emph{Cluster anisotropy.}
The diameter only captures the maximum pairwise distance, but ignores the shape
of the cluster. As a result, clusters can be anisotropic (e.g., elongated or
skewed).

Even if most points are close to the centroid, an elongated cluster can still
contain points that are far apart. Along the short axis, points are very
similar, while along the long axis, points may differ significantly. Thus, two
points in the same cluster can have very different attribute distributions
despite being assigned the same \(P_c(a)\).

\emph{Alignment with distribution gradients.}
Think of \(P(a \mid x)\) as assigning, to every point \(x\) in the embedding
space, a full distribution over attribute values:
\[
x \mapsto P(a \mid x).
\]
If we fix a location \(x\) and look at vectors near \(x\), the
attribute values of those vectors follow a distribution \(P(a \mid x)\). Moving
to a different point \(x'\) may change this distribution. Thus, as we move
through the embedding space, we can think of the attribute distribution as
evolving over space.

\textit{Case 1: Good alignment (low error).}
If a cluster lies within a region where this distribution changes little, then
for any two points \(x, x' \in R_c\), we have \(P(a \mid x) \approx P(a \mid x')\).
In this case, all points in the cluster exhibit similar attribute behavior, and
using a single distribution \(P_c(a)\) is a good approximation.

\textit{Case 2: Poor alignment (high error).}
If a cluster spans a region where the distribution changes significantly, then
points within the cluster may have very different attribute distributions.
However, the model still assigns a single \(P_c(a)\), which forces these
different behaviors to be averaged together and introduces systematic bias.

For example, suppose that in one region of the space the attributes are
concentrated around 10, while in another region they are concentrated around
100. If a cluster spans both regions, then \(P_c(a)\) mixes these two distinct
distributions and fails to accurately represent either.

\emph{Local curvature.}
Finally, the smoothness assumption only captures first-order (linear) variation.
If \(P(a \mid x)\) varies nonlinearly, the Lipschitz bound may underestimate the
true variation within the cluster. As a result, even clusters with small
diameter can incur non-negligible approximation error.

\subsubsection{Estimation error (Statistical Estimation).}
Recall the ideal (population-level) definition of the cluster distribution:
\[
P_c(a) = \mathbb{E}_{x' \in R_c}[P(a \mid x')].
\]
This represents the true cluster-level distribution, which is not directly
observable.

In practice, we estimate this distribution empirically using a histogram over
the points in the cluster:
\[
\hat{P}_c(a) = \frac{1}{n_c} \sum_{i \in C'_c} \mathbf{1}[a_i \in \text{bin}(a)],
\]
where \(n_c = |C'_c|\) is the number of samples in the cluster. This estimator
relies on discretizing the attribute values into bins.

The estimation error is defined as
\[
\varepsilon_{\mathrm{est}}(c)
=
D(\hat{P}_c(a), P_c(a)),
\]
which measures the deviation between the empirical estimate and the true
cluster-level distribution. This error arises from several sources:

\emph{Finite sampling noise (variance).}
We estimate the distribution using a finite number of samples, which introduces
sampling noise that decreases with the sample size:
\[
\varepsilon_{\mathrm{est}}(c) \sim \mathcal{O}\!\left(\frac{1}{\sqrt{n_c}}\right).
\]
This scaling corresponds to the standard deviation of empirical estimates. Thus, large clusters yield more
stable estimates, while small clusters produce noisy distributions.

\emph{Histogram discretization (bias).}
We do not estimate the full distribution, but a binned version
\(P_c(a) \approx P_c(\text{bin}(a))\).
As a result, fine structure within bins is lost and values near bin boundaries
may be distorted. This introduces systematic bias, as the resolution of the
distribution is reduced.

\emph{Interaction with cluster size.}
Small clusters have few samples (\(n_c\)), leading to high variance and unstable
estimates of \(P_c(a)\). Large clusters have more samples and therefore lower
variance, but tend to cover more heterogeneous regions, which increases the
approximation error.

\subsection{Bias--variance trade-off}

The proposed method introduces a trade-off between approximation
(quantization) error and estimation error. This trade-off is controlled by two
parameters: the number of clusters \(k\) and the radius expansion parameter
\(\lambda\).

\subsubsection{Effect of the number of clusters \(k\).}

The number of clusters \(k\) determines the granularity of the partition of the
embedding space.

Increasing \(k\) reduces the cluster diameter, leading to lower quantization
error, but also reduces the number of samples per cluster, resulting in higher
estimation error.

Decreasing \(k\) increases the cluster diameter, leading to higher
quantization error, but increases the number of samples per cluster, resulting
in lower estimation error.

\subsubsection{Effect of radius expansion \(\lambda\).}\label{subsec:radius-expansion}

In addition to clustering, we can expand each cluster region using a parameter
\(\lambda > 1\). We define the expanded region
\[
r'_c = \lambda \, r_c, \quad
R'_c = \{ x : \|x - c_c\| \le r'_c \}.
\]

The diameter of the expanded region satisfies
\[
\mathrm{diam}(R'_c) = \sup_{x,x' \in R'_c} \|x - x'\|.
\]
Using the triangle inequality,
\[
\|x - x'\| \le \|x - c_c\| + \|x' - c_c\| \le r'_c + r'_c = 2r'_c,
\]
which implies
\[
\mathrm{diam}(R'_c) \le 2 r'_c = 2 \lambda r_c.
\]

Thus, \(\lambda\) directly controls an upper bound on the cluster diameter, and
consequently the quantization error.

Increasing \(\lambda\) enlarges the effective region \(R'_c\), increasing the
diameter (and thus the potential quantization error), while also increasing the
number of samples \(n_c\), which reduces estimation variance. Decreasing
\(\lambda\) has the opposite effect: it reduces the diameter and quantization
error, but also reduces the sample size, increasing estimation variance.

Overall, \(\lambda\) acts as a regularization parameter, trading off geometric
fidelity against statistical stability by controlling the spatial extent over
which the conditional distribution is assumed constant.